\section{高宗时期可跳转事件节点}
以下锚点覆盖雍正末至乾隆禅位的财政、制度、边疆、工程、宗教与地方社会事件。来源说明用于提示检索路径，并保留材料对数字、执行与地方反应的限制。
\subsection{1720—1729}
\eventanchor{EQ-1723-YONGZHENG-ACCESSION}{雍正帝即位并开始重整康熙晚年的财…}雍正元年（1723年），雍正帝即位并开始重整康熙晚年的财政、宗室与官僚议题。核心取证：《清世宗实录》雍正元年相关条；宫中朱批奏折、内阁大库题本与《清史稿》相关志传。此类文件确认机构作出或收到的行政动作，不能跳过地方执行与反馈。来源保留：以雍正元年（1723年）的同期文书为定位起点；官修记录、奏报或外交文本服务特定机构立场，具体日期、规模、地方执行与对手视角须以独立材料复核。
\eventanchor{EQ-1723-JI-ZENGYUN-CASE}{年羹尧在西北军政中的权力与朝廷猜…}雍正元年（1723年），年羹尧在西北军政中的权力与朝廷猜忌加剧。核心取证：《清世宗实录》雍正元年相关条；宫中朱批奏折、内阁大库题本与《清史稿》相关志传。此类文件确认机构作出或收到的行政动作，不能跳过地方执行与反馈。来源保留：以雍正元年（1723年）的同期文书为定位起点；官修记录、奏报或外交文本服务特定机构立场，具体日期、规模、地方执行与对手视角须以独立材料复核。
\eventanchor{EQ-1724-CATHOLIC-PROSCRIPTION}{清廷扩大对天主教活动的限制，将传…}雍正二年（1724年），清廷扩大对天主教活动的限制，将传教、地方教民与外来人员置入治理问题。核心取证：《清世宗实录》雍正二年相关条；地方志、督抚奏折及地方档案。编年证词定位国家叙述中的时间，却需同地方或对方材料校验范围。来源保留：以雍正二年（1724年）的同期文书为定位起点；官修记录、奏报或外交文本服务特定机构立场，具体日期、规模、地方执行与对手视角须以独立材料复核。
\eventanchor{EQ-1725-NIAN-GENGYAO-PUNISHMENT}{年羹尧被削职、赐死，其党羽与相关…}雍正三年（1725年），年羹尧被削职、赐死，其党羽与相关官员受到清理。核心取证：《清世宗实录》雍正三年相关条；宫中朱批奏折、内阁大库题本与《清史稿》相关志传。此类文件确认机构作出或收到的行政动作，不能跳过地方执行与反馈。来源保留：以雍正三年（1725年）的同期文书为定位起点；官修记录、奏报或外交文本服务特定机构立场，具体日期、规模、地方执行与对手视角须以独立材料复核。
\eventanchor{EQ-1725-TAN-TEBI-CASE}{隆科多失势并被幽禁，雍正初年近臣…}雍正三年（1725年），隆科多失势并被幽禁，雍正初年近臣权力结构继续调整。核心取证：《清世宗实录》雍正三年相关条；宫中朱批奏折、内阁大库题本与《清史稿》相关志传。此类文件确认机构作出或收到的行政动作，不能跳过地方执行与反馈。来源保留：以雍正三年（1725年）的同期文书为定位起点；官修记录、奏报或外交文本服务特定机构立场，具体日期、规模、地方执行与对手视角须以独立材料复核。
\eventanchor{EQ-1726-HAOXIAN-GUIGONG}{雍正政府推进耗羡归公与养廉银安排…}雍正四年（1726年），雍正政府推进耗羡归公与养廉银安排，试图重组地方财政和官员俸给。核心取证：《清世宗实录》雍正四年相关条；户部奏销、盐政、河工或海关档案。此类文件确认机构作出或收到的行政动作，不能跳过地方执行与反馈。来源保留：以雍正四年（1726年）的同期文书为定位起点；官修记录、奏报或外交文本服务特定机构立场，具体日期、规模、地方执行与对手视角须以独立材料复核。
\eventanchor{EQ-1726-Gaitu-Guiliu-SOUTHWEST}{清廷在西南继续改土归流，将部分土…}雍正四年（1726年），清廷在西南继续改土归流，将部分土司辖地改置流官体系。核心取证：《清世宗实录》雍正四年相关条；军机处档、理藩院档或相关方略。编年证词定位国家叙述中的时间，却需同地方或对方材料校验范围。来源保留：以雍正四年（1726年）的同期文书为定位起点；官修记录、奏报或外交文本服务特定机构立场，具体日期、规模、地方执行与对手视角须以独立材料复核。
\eventanchor{EQ-1727-KYAKHTA-TREATY}{《恰克图条约》确定清俄北部边界与…}雍正五年（1727年），《恰克图条约》确定清俄北部边界与贸易、使团安排。核心取证：《清世宗实录》雍正五年相关条；《筹办夷务始末》相关年卷与中外照会、条约文本。此类文件确认机构作出或收到的行政动作，不能跳过地方执行与反馈。来源保留：以雍正五年（1727年）的同期文书为定位起点；官修记录、奏报或外交文本服务特定机构立场，具体日期、规模、地方执行与对手视角须以独立材料复核。
\eventanchor{EQ-1727-BANNER-GENEALOGIES}{清廷命令编修八旗谱牒，试图界定旗…}雍正五年（1727年），清廷命令编修八旗谱牒，试图界定旗籍、世系与待遇资格。核心取证：《清世宗实录》雍正五年相关条；地方志、督抚奏折及地方档案。编年证词定位国家叙述中的时间，却需同地方或对方材料校验范围。来源保留：以雍正五年（1727年）的同期文书为定位起点；官修记录、奏报或外交文本服务特定机构立场，具体日期、规模、地方执行与对手视角须以独立材料复核。
\eventanchor{EQ-1728-Dzungar-EXPEDITION}{清军在西北对准噶尔展开远征，后勤…}雍正六年（1728年），清军在西北对准噶尔展开远征，后勤与高原草原路线成为军事难题。核心取证：《清世宗实录》雍正六年相关条；军机处档、战事方略及地方奏报。编年证词定位国家叙述中的时间，却需同地方或对方材料校验范围。来源保留：以雍正六年（1728年）的同期文书为定位起点；官修记录、奏报或外交文本服务特定机构立场，具体日期、规模、地方执行与对手视角须以独立材料复核。
\eventanchor{EQ-1728-ZHUNGHAR-TIBET-SETTLEMENT}{清廷在准噶尔威胁下调整驻藏与青海…}雍正六年（1728年），清廷在准噶尔威胁下调整驻藏与青海蒙古的军事部署。核心取证：《清世宗实录》雍正六年相关条；军机处档、理藩院档或相关方略。编年证词定位国家叙述中的时间，却需同地方或对方材料校验范围。来源保留：以雍正六年（1728年）的同期文书为定位起点；官修记录、奏报或外交文本服务特定机构立场，具体日期、规模、地方执行与对手视角须以独立材料复核。
\eventanchor{EQ-1729-JUNJICHU-FOUNDING}{雍正帝设军机房，后来发展为军机处…}雍正七年（1729年），雍正帝设军机房，后来发展为军机处，处理紧急军政文书。核心取证：《清世宗实录》雍正七年相关条；宫中朱批奏折、内阁大库题本与《清史稿》相关志传。此类文件确认机构作出或收到的行政动作，不能跳过地方执行与反馈。来源保留：以雍正七年（1729年）的同期文书为定位起点；官修记录、奏报或外交文本服务特定机构立场，具体日期、规模、地方执行与对手视角须以独立材料复核。
\subsection{1730—1739}
\eventanchor{EQ-1730-SOUTHWEST-MIAO-GOVERNANCE}{清廷围绕苗疆设置、屯兵与地方纠纷…}雍正八年（1730年），清廷围绕苗疆设置、屯兵与地方纠纷加强治理。核心取证：《清世宗实录》雍正八年相关条；军机处档、理藩院档或相关方略。编年证词定位国家叙述中的时间，却需同地方或对方材料校验范围。来源保留：以雍正八年（1730年）的同期文书为定位起点；官修记录、奏报或外交文本服务特定机构立场，具体日期、规模、地方执行与对手视角须以独立材料复核。
\eventanchor{EQ-1731-Dzungar-WAR-REVERSE}{清军在对准噶尔战争中受挫，西北远…}雍正九年（1731年），清军在对准噶尔战争中受挫，西北远征暴露补给与指挥困难。核心取证：《清世宗实录》雍正九年相关条；军机处档、战事方略及地方奏报。编年证词定位国家叙述中的时间，却需同地方或对方材料校验范围。来源保留：以雍正九年（1731年）的同期文书为定位起点；官修记录、奏报或外交文本服务特定机构立场，具体日期、规模、地方执行与对手视角须以独立材料复核。
\eventanchor{EQ-1732-KOBDO-BATTLE}{清军在和通泊等战事失利，西北战局…}雍正十年（1732年），清军在和通泊等战事失利，西北战局一度紧张。核心取证：《清世宗实录》雍正十年相关条；军机处档、战事方略及地方奏报。编年证词定位国家叙述中的时间，却需同地方或对方材料校验范围。来源保留：以雍正十年（1732年）的同期文书为定位起点；官修记录、奏报或外交文本服务特定机构立场，具体日期、规模、地方执行与对手视角须以独立材料复核。
\eventanchor{EQ-1732-AMIN-KHOJA-FLIGHT}{阿敏和卓等人在准噶尔压力下向清境…}雍正十年（1732年），阿敏和卓等人在准噶尔压力下向清境迁移，绿洲政治进入清准竞争。核心取证：《清世宗实录》雍正十年相关条；军机处档、理藩院档或相关方略。编年证词定位国家叙述中的时间，却需同地方或对方材料校验范围。来源保留：以雍正十年（1732年）的同期文书为定位起点；官修记录、奏报或外交文本服务特定机构立场，具体日期、规模、地方执行与对手视角须以独立材料复核。
\eventanchor{EQ-1733-YONGZHENG-FISCAL-INSPECTION}{雍正政府持续以奏销、亏空追查和官…}雍正十一年（1733年），雍正政府持续以奏销、亏空追查和官员考成为手段整顿地方财政。核心取证：《清世宗实录》雍正十一年相关条；户部奏销、盐政、河工或海关档案。此类文件确认机构作出或收到的行政动作，不能跳过地方执行与反馈。来源保留：以雍正十一年（1733年）的同期文书为定位起点；官修记录、奏报或外交文本服务特定机构立场，具体日期、规模、地方执行与对手视角须以独立材料复核。
\eventanchor{EQ-1734-CHRISTIAN-ENFORCEMENT}{地方官继续执行限制天主教的政策，…}雍正十二年（1734年），地方官继续执行限制天主教的政策，传教士和教民的活动受地域差异影响。核心取证：《清世宗实录》雍正十二年相关条；地方志、督抚奏折及地方档案。编年证词定位国家叙述中的时间，却需同地方或对方材料校验范围。来源保留：以雍正十二年（1734年）的同期文书为定位起点；官修记录、奏报或外交文本服务特定机构立场，具体日期、规模、地方执行与对手视角须以独立材料复核。
\eventanchor{EQ-1735-QIANLONG-ACCESSION}{乾隆帝即位，雍正时期形成的财政与…}雍正十三年（1735年），乾隆帝即位，雍正时期形成的财政与军机决策机制进入新朝。核心取证：《清世宗实录》雍正十三年相关条；宫中朱批奏折、内阁大库题本与《清史稿》相关志传。此类文件确认机构作出或收到的行政动作，不能跳过地方执行与反馈。来源保留：以雍正十三年（1735年）的同期文书为定位起点；官修记录、奏报或外交文本服务特定机构立场，具体日期、规模、地方执行与对手视角须以独立材料复核。
\eventanchor{EQ-1735-MIAO-UPRISING-RONGJIANG}{贵州荣江等地苗侗社区反抗清军和地…}雍正十三年（1735年），贵州荣江等地苗侗社区反抗清军和地方治理，镇压造成严重暴力。核心取证：《清世宗实录》雍正十三年相关条；军机处档、战事方略及地方奏报。编年证词定位国家叙述中的时间，却需同地方或对方材料校验范围。来源保留：以雍正十三年（1735年）的同期文书为定位起点；官修记录、奏报或外交文本服务特定机构立场，具体日期、规模、地方执行与对手视角须以独立材料复核。
\eventanchor{EQ-1736-AMNESTY-AND-RECTIFICATION}{乾隆初年颁布恩诏并调整雍正后期政…}乾隆元年（1736年），乾隆初年颁布恩诏并调整雍正后期政治案件的处置。核心取证：《清高宗实录》乾隆元年相关条；宫中朱批奏折、内阁大库题本与《清史稿》相关志传。此类文件确认机构作出或收到的行政动作，不能跳过地方执行与反馈。来源保留：以乾隆元年（1736年）的同期文书为定位起点；官修记录、奏报或外交文本服务特定机构立场，具体日期、规模、地方执行与对手视角须以独立材料复核。
\eventanchor{EQ-1737-Dzungar-NEGOTIATION}{清廷在战争压力下继续处理与准噶尔…}乾隆二年（1737年），清廷在战争压力下继续处理与准噶尔的使节、盟旗和军事情报。核心取证：《清高宗实录》乾隆二年相关条；《筹办夷务始末》相关年卷与中外照会、条约文本。此类文件确认机构作出或收到的行政动作，不能跳过地方执行与反馈。来源保留：以乾隆二年（1737年）的同期文书为定位起点；官修记录、奏报或外交文本服务特定机构立场，具体日期、规模、地方执行与对手视角须以独立材料复核。
\eventanchor{EQ-1738-FIRST-JINCHUAN-WAR}{清军开始大规模干预大金川，山地战…}乾隆三年（1738年），清军开始大规模干预大金川，山地战与土司关系限制军事推进。核心取证：《清高宗实录》乾隆三年相关条；军机处档、战事方略及地方奏报。编年证词定位国家叙述中的时间，却需同地方或对方材料校验范围。来源保留：以乾隆三年（1738年）的同期文书为定位起点；官修记录、奏报或外交文本服务特定机构立场，具体日期、规模、地方执行与对手视角须以独立材料复核。
\eventanchor{EQ-1739-JINCHUAN-SUPPLY}{大金川战事的粮运、道路和军费成为…}乾隆四年（1739年），大金川战事的粮运、道路和军费成为中央持续关注的难题。核心取证：《清高宗实录》乾隆四年相关条；户部奏销、盐政、河工或海关档案。此类文件确认机构作出或收到的行政动作，不能跳过地方执行与反馈。来源保留：以乾隆四年（1739年）的同期文书为定位起点；官修记录、奏报或外交文本服务特定机构立场，具体日期、规模、地方执行与对手视角须以独立材料复核。
\subsection{1740—1749}
\eventanchor{EQ-1740-NORTH-CHINA-FAMINE}{华北旱灾和粮价问题促使地方奏报与…}乾隆五年（1740年），华北旱灾和粮价问题促使地方奏报与赈济措施增加。核心取证：《清高宗实录》乾隆五年相关条；地方志、督抚奏折及地方档案。编年证词定位国家叙述中的时间，却需同地方或对方材料校验范围。来源保留：以乾隆五年（1740年）的同期文书为定位起点；官修记录、奏报或外交文本服务特定机构立场，具体日期、规模、地方执行与对手视角须以独立材料复核。
\eventanchor{EQ-1741-MIAO-FRONTIER-CAMPAIGN}{清廷在西南继续以驻兵、设治和军事…}乾隆六年（1741年），清廷在西南继续以驻兵、设治和军事行动处理苗疆冲突。核心取证：《清高宗实录》乾隆六年相关条；军机处档、理藩院档或相关方略。编年证词定位国家叙述中的时间，却需同地方或对方材料校验范围。来源保留：以乾隆六年（1741年）的同期文书为定位起点；官修记录、奏报或外交文本服务特定机构立场，具体日期、规模、地方执行与对手视角须以独立材料复核。
\eventanchor{EQ-1742-BANNER-EXIT-ORDER}{清廷允许部分汉军旗人出旗，调整旗…}乾隆七年（1742年），清廷允许部分汉军旗人出旗，调整旗籍与生计安排。核心取证：《清高宗实录》乾隆七年相关条；地方志、督抚奏折及地方档案。编年证词定位国家叙述中的时间，却需同地方或对方材料校验范围。来源保留：以乾隆七年（1742年）的同期文书为定位起点；官修记录、奏报或外交文本服务特定机构立场，具体日期、规模、地方执行与对手视角须以独立材料复核。
\eventanchor{EQ-1743-YELLOW-RIVER-WORKS}{河工官员就黄河堤防、河道和经费持…}乾隆八年（1743年），河工官员就黄河堤防、河道和经费持续奏报并实施工程。核心取证：《清高宗实录》乾隆八年相关条；户部奏销、盐政、河工或海关档案。此类文件确认机构作出或收到的行政动作，不能跳过地方执行与反馈。来源保留：以乾隆八年（1743年）的同期文书为定位起点；官修记录、奏报或外交文本服务特定机构立场，具体日期、规模、地方执行与对手视角须以独立材料复核。
\eventanchor{EQ-1744-JESUIT-CALENDAR-SERVICE}{宫廷继续利用部分西洋传教士的历算…}乾隆九年（1744年），宫廷继续利用部分西洋传教士的历算与技术服务，同时维持宗教限制。核心取证：《清高宗实录》乾隆九年相关条；内阁档、文集或报刊相关记录。编年证词定位国家叙述中的时间，却需同地方或对方材料校验范围。来源保留：以乾隆九年（1744年）的同期文书为定位起点；官修记录、奏报或外交文本服务特定机构立场，具体日期、规模、地方执行与对手视角须以独立材料复核。
\eventanchor{EQ-1745-JINCHUAN-NEGOTIATION}{清廷在大金川战局中尝试通过招抚、…}乾隆十年（1745年），清廷在大金川战局中尝试通过招抚、盟约和军事压力重组地方关系。核心取证：《清高宗实录》乾隆十年相关条；军机处档、理藩院档或相关方略。编年证词定位国家叙述中的时间，却需同地方或对方材料校验范围。来源保留：以乾隆十年（1745年）的同期文书为定位起点；官修记录、奏报或外交文本服务特定机构立场，具体日期、规模、地方执行与对手视角须以独立材料复核。
\eventanchor{EQ-1746-FIRST-JINCHUAN-SETTLEMENT}{第一次金川战争以和议和行政安排告…}乾隆十一年（1746年），第一次金川战争以和议和行政安排告一段落。核心取证：《清高宗实录》乾隆十一年相关条；军机处档、战事方略及地方奏报。编年证词定位国家叙述中的时间，却需同地方或对方材料校验范围。来源保留：以乾隆十一年（1746年）的同期文书为定位起点；官修记录、奏报或外交文本服务特定机构立场，具体日期、规模、地方执行与对手视角须以独立材料复核。
\eventanchor{EQ-1747-SOUTHWEST-MINING-REGULATION}{清廷和地方官围绕银铜矿开采、税课…}乾隆十二年（1747年），清廷和地方官围绕银铜矿开采、税课与矿工流动加强管理。核心取证：《清高宗实录》乾隆十二年相关条；户部奏销、盐政、河工或海关档案。此类文件确认机构作出或收到的行政动作，不能跳过地方执行与反馈。来源保留：以乾隆十二年（1747年）的同期文书为定位起点；官修记录、奏报或外交文本服务特定机构立场，具体日期、规模、地方执行与对手视角须以独立材料复核。
\eventanchor{EQ-1748-SONG-AND-CASE}{乾隆朝通过大案处分官员，强化对地…}乾隆十三年（1748年），乾隆朝通过大案处分官员，强化对地方文书、钱粮和吏治的监督。核心取证：《清高宗实录》乾隆十三年相关条；宫中朱批奏折、内阁大库题本与《清史稿》相关志传。此类文件确认机构作出或收到的行政动作，不能跳过地方执行与反馈。来源保留：以乾隆十三年（1748年）的同期文书为定位起点；官修记录、奏报或外交文本服务特定机构立场，具体日期、规模、地方执行与对手视角须以独立材料复核。
\eventanchor{EQ-1749-JIANGNAN-TAX-REVIEW}{江南钱粮、漕运和士绅关系继续在户…}乾隆十四年（1749年），江南钱粮、漕运和士绅关系继续在户部与地方奏折中被讨论。核心取证：《清高宗实录》乾隆十四年相关条；户部奏销、盐政、河工或海关档案。此类文件确认机构作出或收到的行政动作，不能跳过地方执行与反馈。来源保留：以乾隆十四年（1749年）的同期文书为定位起点；官修记录、奏报或外交文本服务特定机构立场，具体日期、规模、地方执行与对手视角须以独立材料复核。
\subsection{1750—1759}
\eventanchor{EQ-1750-LHASA-UPRISING}{拉萨发生政治冲突，驻藏大臣和清军…}乾隆十五年（1750年），拉萨发生政治冲突，驻藏大臣和清军介入处置。核心取证：《清高宗实录》乾隆十五年相关条；军机处档、理藩院档或相关方略。编年证词定位国家叙述中的时间，却需同地方或对方材料校验范围。来源保留：以乾隆十五年（1750年）的同期文书为定位起点；官修记录、奏报或外交文本服务特定机构立场，具体日期、规模、地方执行与对手视角须以独立材料复核。
\eventanchor{EQ-1751-TIBET-ORDINANCE}{清廷在拉萨事变后重整驻藏大臣和噶…}乾隆十六年（1751年），清廷在拉萨事变后重整驻藏大臣和噶厦相关制度安排。核心取证：《清高宗实录》乾隆十六年相关条；军机处档、理藩院档或相关方略。编年证词定位国家叙述中的时间，却需同地方或对方材料校验范围。来源保留：以乾隆十六年（1751年）的同期文书为定位起点；官修记录、奏报或外交文本服务特定机构立场，具体日期、规模、地方执行与对手视角须以独立材料复核。
\eventanchor{EQ-1752-TIBETAN-ADMINISTRATION}{驻藏系统与青海、川藏交通的军政协…}乾隆十七年（1752年），驻藏系统与青海、川藏交通的军政协调继续调整。核心取证：《清高宗实录》乾隆十七年相关条；军机处档、理藩院档或相关方略。编年证词定位国家叙述中的时间，却需同地方或对方材料校验范围。来源保留：以乾隆十七年（1752年）的同期文书为定位起点；官修记录、奏报或外交文本服务特定机构立场，具体日期、规模、地方执行与对手视角须以独立材料复核。
\eventanchor{EQ-1753-SALT-ADMINISTRATION}{盐政、引岸和课额问题在中央与产销…}乾隆十八年（1753年），盐政、引岸和课额问题在中央与产销地区持续受到整顿。核心取证：《清高宗实录》乾隆十八年相关条；户部奏销、盐政、河工或海关档案。此类文件确认机构作出或收到的行政动作，不能跳过地方执行与反馈。来源保留：以乾隆十八年（1753年）的同期文书为定位起点；官修记录、奏报或外交文本服务特定机构立场，具体日期、规模、地方执行与对手视角须以独立材料复核。
\eventanchor{EQ-1754-AMURSANA-DEFECTION}{阿睦尔撒纳等准噶尔贵族向清廷归附…}乾隆十九年（1754年），阿睦尔撒纳等准噶尔贵族向清廷归附，改变西北战争的联盟格局。核心取证：《清高宗实录》乾隆十九年相关条；军机处档、理藩院档或相关方略。编年证词定位国家叙述中的时间，却需同地方或对方材料校验范围。来源保留：以乾隆十九年（1754年）的同期文书为定位起点；官修记录、奏报或外交文本服务特定机构立场，具体日期、规模、地方执行与对手视角须以独立材料复核。
\eventanchor{EQ-1755-ILI-EXPEDITION}{清军分路进入伊犁，击溃达瓦齐政权…}乾隆二十年（1755年），清军分路进入伊犁，击溃达瓦齐政权并占领准噶尔核心地区。核心取证：《清高宗实录》乾隆二十年相关条；军机处档、战事方略及地方奏报。编年证词定位国家叙述中的时间，却需同地方或对方材料校验范围。来源保留：以乾隆二十年（1755年）的同期文书为定位起点；官修记录、奏报或外交文本服务特定机构立场，具体日期、规模、地方执行与对手视角须以独立材料复核。
\eventanchor{EQ-1755-AMURSANA-REVOLT}{阿睦尔撒纳反清，准噶尔战后秩序重…}乾隆二十年（1755年），阿睦尔撒纳反清，准噶尔战后秩序重新陷入军事冲突。核心取证：《清高宗实录》乾隆二十年相关条；军机处档、战事方略及地方奏报。编年证词定位国家叙述中的时间，却需同地方或对方材料校验范围。来源保留：以乾隆二十年（1755年）的同期文书为定位起点；官修记录、奏报或外交文本服务特定机构立场，具体日期、规模、地方执行与对手视角须以独立材料复核。
\eventanchor{EQ-1756-BANNER-CIVILIAN-REGISTRATION}{清廷命令部分八旗副户口改归民籍，…}乾隆二十一年（1756年），清廷命令部分八旗副户口改归民籍，继续缓解旗籍供养压力。核心取证：《清高宗实录》乾隆二十一年相关条；地方志、督抚奏折及地方档案。编年证词定位国家叙述中的时间，却需同地方或对方材料校验范围。来源保留：以乾隆二十一年（1756年）的同期文书为定位起点；官修记录、奏报或外交文本服务特定机构立场，具体日期、规模、地方执行与对手视角须以独立材料复核。
\eventanchor{EQ-1757-CANTON-TRADE-RESTRICTION}{清廷将西洋贸易集中于广州口岸，形…}乾隆二十二年（1757年），清廷将西洋贸易集中于广州口岸，形成后来所谓“一口通商”的制度框架。核心取证：《清高宗实录》乾隆二十二年相关条；《筹办夷务始末》相关年卷与中外照会、条约文本。此类文件确认机构作出或收到的行政动作，不能跳过地方执行与反馈。来源保留：以乾隆二十二年（1757年）的同期文书为定位起点；官修记录、奏报或外交文本服务特定机构立场，具体日期、规模、地方执行与对手视角须以独立材料复核。
\eventanchor{EQ-1757-AMURSANA-DEATH}{阿睦尔撒纳逃入俄境后死亡，清准战…}乾隆二十二年（1757年），阿睦尔撒纳逃入俄境后死亡，清准战争的关键对手退出。核心取证：《清高宗实录》乾隆二十二年相关条；军机处档、理藩院档或相关方略。编年证词定位国家叙述中的时间，却需同地方或对方材料校验范围。来源保留：以乾隆二十二年（1757年）的同期文书为定位起点；官修记录、奏报或外交文本服务特定机构立场，具体日期、规模、地方执行与对手视角须以独立材料复核。
\eventanchor{EQ-1758-KHOJA-REBELLION}{大小和卓反清，清军由北路向天山南…}乾隆二十三年（1758年），大小和卓反清，清军由北路向天山南路推进。核心取证：《清高宗实录》乾隆二十三年相关条；军机处档、战事方略及地方奏报。编年证词定位国家叙述中的时间，却需同地方或对方材料校验范围。来源保留：以乾隆二十三年（1758年）的同期文书为定位起点；官修记录、奏报或外交文本服务特定机构立场，具体日期、规模、地方执行与对手视角须以独立材料复核。
\eventanchor{EQ-1759-KASHGAR-YARKAND-CAMPAIGN}{清军攻入喀什噶尔、叶尔羌等地，和…}乾隆二十四年（1759年），清军攻入喀什噶尔、叶尔羌等地，和卓政权覆亡。核心取证：《清高宗实录》乾隆二十四年相关条；军机处档、战事方略及地方奏报。编年证词定位国家叙述中的时间，却需同地方或对方材料校验范围。来源保留：以乾隆二十四年（1759年）的同期文书为定位起点；官修记录、奏报或外交文本服务特定机构立场，具体日期、规模、地方执行与对手视角须以独立材料复核。
\subsection{1760—1769}
\eventanchor{EQ-1760-ILI-GENERAL}{清廷设伊犁将军并展开驻防、屯田与…}乾隆二十五年（1760年），清廷设伊犁将军并展开驻防、屯田与移民安置。核心取证：《清高宗实录》乾隆二十五年相关条；军机处档、理藩院档或相关方略。编年证词定位国家叙述中的时间，却需同地方或对方材料校验范围。来源保留：以乾隆二十五年（1760年）的同期文书为定位起点；官修记录、奏报或外交文本服务特定机构立场，具体日期、规模、地方执行与对手视角须以独立材料复核。
\eventanchor{EQ-1761-XINJIANG-GARRISONS}{清廷在新疆配置满汉蒙军队和绿营驻…}乾隆二十六年（1761年），清廷在新疆配置满汉蒙军队和绿营驻防，建立军政节点。核心取证：《清高宗实录》乾隆二十六年相关条；军机处档、理藩院档或相关方略。编年证词定位国家叙述中的时间，却需同地方或对方材料校验范围。来源保留：以乾隆二十六年（1761年）的同期文书为定位起点；官修记录、奏报或外交文本服务特定机构立场，具体日期、规模、地方执行与对手视角须以独立材料复核。
\eventanchor{EQ-1762-URUMQI-FOUNDATION}{迪化等北疆城市的军政建设加快，成…}乾隆二十七年（1762年），迪化等北疆城市的军政建设加快，成为屯田和交通枢纽。核心取证：《清高宗实录》乾隆二十七年相关条；军机处档、理藩院档或相关方略。编年证词定位国家叙述中的时间，却需同地方或对方材料校验范围。来源保留：以乾隆二十七年（1762年）的同期文书为定位起点；官修记录、奏报或外交文本服务特定机构立场，具体日期、规模、地方执行与对手视角须以独立材料复核。
\eventanchor{EQ-1763-XINJIANG-SETTLEMENT}{清廷继续在新疆安置屯田人口、军户…}乾隆二十八年（1763年），清廷继续在新疆安置屯田人口、军户和商业网络。核心取证：《清高宗实录》乾隆二十八年相关条；地方志、督抚奏折及地方档案。编年证词定位国家叙述中的时间，却需同地方或对方材料校验范围。来源保留：以乾隆二十八年（1763年）的同期文书为定位起点；官修记录、奏报或外交文本服务特定机构立场，具体日期、规模、地方执行与对手视角须以独立材料复核。
\eventanchor{EQ-1764-MANCHU-BANNER-RESETTLEMENT}{部分八旗兵丁被调往西北驻防，旗籍…}乾隆二十九年（1764年），部分八旗兵丁被调往西北驻防，旗籍家庭的迁移与供给成为边务问题。核心取证：《清高宗实录》乾隆二十九年相关条；地方志、督抚奏折及地方档案。编年证词定位国家叙述中的时间，却需同地方或对方材料校验范围。来源保留：以乾隆二十九年（1764年）的同期文书为定位起点；官修记录、奏报或外交文本服务特定机构立场，具体日期、规模、地方执行与对手视角须以独立材料复核。
\eventanchor{EQ-1765-USH-REBELLION}{乌什爆发反清事件，清军镇压并重组…}乾隆三十年（1765年），乌什爆发反清事件，清军镇压并重组当地治理。核心取证：《清高宗实录》乾隆三十年相关条；军机处档、战事方略及地方奏报。编年证词定位国家叙述中的时间，却需同地方或对方材料校验范围。来源保留：以乾隆三十年（1765年）的同期文书为定位起点；官修记录、奏报或外交文本服务特定机构立场，具体日期、规模、地方执行与对手视角须以独立材料复核。
\eventanchor{EQ-1766-BURMA-WAR-BEGIN}{清军对缅甸发动远征，西南边境战争…}乾隆三十一年（1766年），清军对缅甸发动远征，西南边境战争展开。核心取证：《清高宗实录》乾隆三十一年相关条；军机处档、战事方略及地方奏报。编年证词定位国家叙述中的时间，却需同地方或对方材料校验范围。来源保留：以乾隆三十一年（1766年）的同期文书为定位起点；官修记录、奏报或外交文本服务特定机构立场，具体日期、规模、地方执行与对手视角须以独立材料复核。
\eventanchor{EQ-1767-BURMA-CAMPAIGN}{清缅战争持续，清军在边境和山地路…}乾隆三十二年（1767年），清缅战争持续，清军在边境和山地路线中遭遇重大困难。核心取证：《清高宗实录》乾隆三十二年相关条；军机处档、战事方略及地方奏报。编年证词定位国家叙述中的时间，却需同地方或对方材料校验范围。来源保留：以乾隆三十二年（1767年）的同期文书为定位起点；官修记录、奏报或外交文本服务特定机构立场，具体日期、规模、地方执行与对手视角须以独立材料复核。
\eventanchor{EQ-1768-LITERARY-INQUISITION}{乾隆朝以文字狱案件强化对著述、地…}乾隆三十三年（1768年），乾隆朝以文字狱案件强化对著述、地方士人和政治语言的控制。核心取证：《清高宗实录》乾隆三十三年相关条；内阁档、文集或报刊相关记录。编年证词定位国家叙述中的时间，却需同地方或对方材料校验范围。来源保留：以乾隆三十三年（1768年）的同期文书为定位起点；官修记录、奏报或外交文本服务特定机构立场，具体日期、规模、地方执行与对手视角须以独立材料复核。
\eventanchor{EQ-1768-SECOND-JINCHUAN-WAR}{第二次金川战争开始，清廷投入更大…}乾隆三十三年（1768年），第二次金川战争开始，清廷投入更大兵力和资源。核心取证：《清高宗实录》乾隆三十三年相关条；军机处档、战事方略及地方奏报。编年证词定位国家叙述中的时间，却需同地方或对方材料校验范围。来源保留：以乾隆三十三年（1768年）的同期文书为定位起点；官修记录、奏报或外交文本服务特定机构立场，具体日期、规模、地方执行与对手视角须以独立材料复核。
\eventanchor{EQ-1769-BURMA-SETTLEMENT}{清缅战争在反复交战后趋向停战和使…}乾隆三十四年（1769年），清缅战争在反复交战后趋向停战和使节往来安排。核心取证：《清高宗实录》乾隆三十四年相关条；《筹办夷务始末》相关年卷与中外照会、条约文本。此类文件确认机构作出或收到的行政动作，不能跳过地方执行与反馈。来源保留：以乾隆三十四年（1769年）的同期文书为定位起点；官修记录、奏报或外交文本服务特定机构立场，具体日期、规模、地方执行与对手视角须以独立材料复核。
\eventanchor{EQ-1769-JINCHUAN-ESCALATION}{金川战场的道路、炮兵、粮运与将领…}乾隆三十四年（1769年），金川战场的道路、炮兵、粮运与将领调度不断升级。核心取证：《清高宗实录》乾隆三十四年相关条；军机处档、战事方略及地方奏报。编年证词定位国家叙述中的时间，却需同地方或对方材料校验范围。来源保留：以乾隆三十四年（1769年）的同期文书为定位起点；官修记录、奏报或外交文本服务特定机构立场，具体日期、规模、地方执行与对手视角须以独立材料复核。
\subsection{1770—1779}
\eventanchor{EQ-1770-SICHUAN-SUPPLY-BURDEN}{金川军需加重四川及邻省的转运、民…}乾隆三十五年（1770年），金川军需加重四川及邻省的转运、民夫和财政负担。核心取证：《清高宗实录》乾隆三十五年相关条；户部奏销、盐政、河工或海关档案。此类文件确认机构作出或收到的行政动作，不能跳过地方执行与反馈。来源保留：以乾隆三十五年（1770年）的同期文书为定位起点；官修记录、奏报或外交文本服务特定机构立场，具体日期、规模、地方执行与对手视角须以独立材料复核。
\eventanchor{EQ-1771-TORGHUT-RETURN}{土尔扈特部自伏尔加回归，清廷在伊…}乾隆三十六年（1771年），土尔扈特部自伏尔加回归，清廷在伊犁等地安置部众。核心取证：《清高宗实录》乾隆三十六年相关条；军机处档、理藩院档或相关方略。编年证词定位国家叙述中的时间，却需同地方或对方材料校验范围。来源保留：以乾隆三十六年（1771年）的同期文书为定位起点；官修记录、奏报或外交文本服务特定机构立场，具体日期、规模、地方执行与对手视角须以独立材料复核。
\eventanchor{EQ-1772-SIKU-QUANSHU-ORDER}{乾隆帝下令征集、编纂《四库全书》…}乾隆三十七年（1772年），乾隆帝下令征集、编纂《四库全书》，建立全国性的图书征集机制。核心取证：《清高宗实录》乾隆三十七年相关条；内阁档、文集或报刊相关记录。编年证词定位国家叙述中的时间，却需同地方或对方材料校验范围。来源保留：以乾隆三十七年（1772年）的同期文书为定位起点；官修记录、奏报或外交文本服务特定机构立场，具体日期、规模、地方执行与对手视角须以独立材料复核。
\eventanchor{EQ-1773-SICHUAN-RECRUITMENT}{金川战事促使清廷继续在四川及邻省…}乾隆三十八年（1773年），金川战事促使清廷继续在四川及邻省征调兵丁和夫役。核心取证：《清高宗实录》乾隆三十八年相关条；军机处档、战事方略及地方奏报。编年证词定位国家叙述中的时间，却需同地方或对方材料校验范围。来源保留：以乾隆三十八年（1773年）的同期文书为定位起点；官修记录、奏报或外交文本服务特定机构立场，具体日期、规模、地方执行与对手视角须以独立材料复核。
\eventanchor{EQ-1774-WANG-LUN-UPRISING}{王伦在山东起事，清廷迅速调兵镇压}乾隆三十九年（1774年），王伦在山东起事，清廷迅速调兵镇压。核心取证：《清高宗实录》乾隆三十九年相关条；地方志、督抚奏折及地方档案。编年证词定位国家叙述中的时间，却需同地方或对方材料校验范围。来源保留：以乾隆三十九年（1774年）的同期文书为定位起点；官修记录、奏报或外交文本服务特定机构立场，具体日期、规模、地方执行与对手视角须以独立材料复核。
\eventanchor{EQ-1775-SIKU-COLLECTION}{《四库全书》征书与审查在各省展开…}乾隆四十年（1775年），《四库全书》征书与审查在各省展开，地方藏书和文人网络受到影响。核心取证：《清高宗实录》乾隆四十年相关条；内阁档、文集或报刊相关记录。编年证词定位国家叙述中的时间，却需同地方或对方材料校验范围。来源保留：以乾隆四十年（1775年）的同期文书为定位起点；官修记录、奏报或外交文本服务特定机构立场，具体日期、规模、地方执行与对手视角须以独立材料复核。
\eventanchor{EQ-1776-JINCHUAN-CONQUEST}{清军攻克金川，第二次金川战争结束}乾隆四十一年（1776年），清军攻克金川，第二次金川战争结束。核心取证：《清高宗实录》乾隆四十一年相关条；军机处档、战事方略及地方奏报。编年证词定位国家叙述中的时间，却需同地方或对方材料校验范围。来源保留：以乾隆四十一年（1776年）的同期文书为定位起点；官修记录、奏报或外交文本服务特定机构立场，具体日期、规模、地方执行与对手视角须以独立材料复核。
\eventanchor{EQ-1777-JINCHUAN-ADMINISTRATIVE-RESET}{清廷在金川地区推行改土归流、设治…}乾隆四十二年（1777年），清廷在金川地区推行改土归流、设治和驻防安排。核心取证：《清高宗实录》乾隆四十二年相关条；军机处档、理藩院档或相关方略。编年证词定位国家叙述中的时间，却需同地方或对方材料校验范围。来源保留：以乾隆四十二年（1777年）的同期文书为定位起点；官修记录、奏报或外交文本服务特定机构立场，具体日期、规模、地方执行与对手视角须以独立材料复核。
\eventanchor{EQ-1778-SICHUAN-TAX-RECOVERY}{金川战后四川财政和民力恢复成为地…}乾隆四十三年（1778年），金川战后四川财政和民力恢复成为地方官奏报的重要议题。核心取证：《清高宗实录》乾隆四十三年相关条；户部奏销、盐政、河工或海关档案。此类文件确认机构作出或收到的行政动作，不能跳过地方执行与反馈。来源保留：以乾隆四十三年（1778年）的同期文书为定位起点；官修记录、奏报或外交文本服务特定机构立场，具体日期、规模、地方执行与对手视角须以独立材料复核。
\eventanchor{EQ-1779-VIETNAM-INTERVENTION}{清廷介入安南政局并派兵南下，孙士…}乾隆四十四年（1779年），清廷介入安南政局并派兵南下，孙士毅军一度进入升龙。核心取证：《清高宗实录》乾隆四十四年相关条；军机处档、战事方略及地方奏报。编年证词定位国家叙述中的时间，却需同地方或对方材料校验范围。来源保留：以乾隆四十四年（1779年）的同期文书为定位起点；官修记录、奏报或外交文本服务特定机构立场，具体日期、规模、地方执行与对手视角须以独立材料复核。
\subsection{1780—1789}
\eventanchor{EQ-1780-VIETNAM-WITHDRAWAL}{清军撤离安南，双方通过使节和册封…}乾隆四十五年（1780年），清军撤离安南，双方通过使节和册封语言恢复关系。核心取证：《清高宗实录》乾隆四十五年相关条；《筹办夷务始末》相关年卷与中外照会、条约文本。此类文件确认机构作出或收到的行政动作，不能跳过地方执行与反馈。来源保留：以乾隆四十五年（1780年）的同期文书为定位起点；官修记录、奏报或外交文本服务特定机构立场，具体日期、规模、地方执行与对手视角须以独立材料复核。
\eventanchor{EQ-1781-GANSU-MUSLIM-UPRISING}{甘肃爆发回民起事，清军实施镇压与…}乾隆四十六年（1781年），甘肃爆发回民起事，清军实施镇压与后续治理。核心取证：《清高宗实录》乾隆四十六年相关条；军机处档、战事方略及地方奏报。编年证词定位国家叙述中的时间，却需同地方或对方材料校验范围。来源保留：以乾隆四十六年（1781年）的同期文书为定位起点；官修记录、奏报或外交文本服务特定机构立场，具体日期、规模、地方执行与对手视角须以独立材料复核。
\eventanchor{EQ-1782-GANSU-POSTWAR-GOVERNANCE}{清廷在甘肃起事后调整驻兵、司法和…}乾隆四十七年（1782年），清廷在甘肃起事后调整驻兵、司法和地方控制措施。核心取证：《清高宗实录》乾隆四十七年相关条；地方志、督抚奏折及地方档案。编年证词定位国家叙述中的时间，却需同地方或对方材料校验范围。来源保留：以乾隆四十七年（1782年）的同期文书为定位起点；官修记录、奏报或外交文本服务特定机构立场，具体日期、规模、地方执行与对手视角须以独立材料复核。
\eventanchor{EQ-1783-SALT-MONOPOLY-REVIEW}{清廷继续讨论盐政弊端、商人负担和…}乾隆四十八年（1783年），清廷继续讨论盐政弊端、商人负担和课额征解。核心取证：《清高宗实录》乾隆四十八年相关条；户部奏销、盐政、河工或海关档案。此类文件确认机构作出或收到的行政动作，不能跳过地方执行与反馈。来源保留：以乾隆四十八年（1783年）的同期文书为定位起点；官修记录、奏报或外交文本服务特定机构立场，具体日期、规模、地方执行与对手视角须以独立材料复核。
\eventanchor{EQ-1785-TAIWAN-GARRISON}{清廷调整台湾驻防、渡台与地方治安…}乾隆五十年（1785年），清廷调整台湾驻防、渡台与地方治安安排。核心取证：《清高宗实录》乾隆五十年相关条；军机处档、理藩院档或相关方略。编年证词定位国家叙述中的时间，却需同地方或对方材料校验范围。来源保留：以乾隆五十年（1785年）的同期文书为定位起点；官修记录、奏报或外交文本服务特定机构立场，具体日期、规模、地方执行与对手视角须以独立材料复核。
\eventanchor{EQ-1786-LIN-SHUANGWEN-UPRISING}{林爽文起事在台湾爆发并迅速扩展}乾隆五十一年（1786年），林爽文起事在台湾爆发并迅速扩展。核心取证：《清高宗实录》乾隆五十一年相关条；军机处档、战事方略及地方奏报。编年证词定位国家叙述中的时间，却需同地方或对方材料校验范围。来源保留：以乾隆五十一年（1786年）的同期文书为定位起点；官修记录、奏报或外交文本服务特定机构立场，具体日期、规模、地方执行与对手视角须以独立材料复核。
\eventanchor{EQ-1786-FIRST-GURKHA-WAR}{廓尔喀军进入西藏边境，清廷开始处…}乾隆五十一年（1786年），廓尔喀军进入西藏边境，清廷开始处理高原军事危机。核心取证：《清高宗实录》乾隆五十一年相关条；军机处档、理藩院档或相关方略。编年证词定位国家叙述中的时间，却需同地方或对方材料校验范围。来源保留：以乾隆五十一年（1786年）的同期文书为定位起点；官修记录、奏报或外交文本服务特定机构立场，具体日期、规模、地方执行与对手视角须以独立材料复核。
\eventanchor{EQ-1787-TAIWAN-CAMPAIGN}{福康安等率军赴台镇压林爽文起事，…}乾隆五十二年（1787年），福康安等率军赴台镇压林爽文起事，海峡运输成为关键。核心取证：《清高宗实录》乾隆五十二年相关条；军机处档、战事方略及地方奏报。编年证词定位国家叙述中的时间，却需同地方或对方材料校验范围。来源保留：以乾隆五十二年（1787年）的同期文书为定位起点；官修记录、奏报或外交文本服务特定机构立场，具体日期、规模、地方执行与对手视角须以独立材料复核。
\eventanchor{EQ-1788-LIN-SHUANGWEN-DEFEAT}{清军击败林爽文，台湾起事基本结束}乾隆五十三年（1788年），清军击败林爽文，台湾起事基本结束。核心取证：《清高宗实录》乾隆五十三年相关条；军机处档、战事方略及地方奏报。编年证词定位国家叙述中的时间，却需同地方或对方材料校验范围。来源保留：以乾隆五十三年（1788年）的同期文书为定位起点；官修记录、奏报或外交文本服务特定机构立场，具体日期、规模、地方执行与对手视角须以独立材料复核。
\eventanchor{EQ-1788-SECOND-GURKHA-WAR}{廓尔喀再次入侵西藏，清军组织更大…}乾隆五十三年（1788年），廓尔喀再次入侵西藏，清军组织更大规模高原远征。核心取证：《清高宗实录》乾隆五十三年相关条；军机处档、战事方略及地方奏报。编年证词定位国家叙述中的时间，却需同地方或对方材料校验范围。来源保留：以乾隆五十三年（1788年）的同期文书为定位起点；官修记录、奏报或外交文本服务特定机构立场，具体日期、规模、地方执行与对手视角须以独立材料复核。
\eventanchor{EQ-1789-TIBET-NEPAL-CAMPAIGN}{清军越过喜马拉雅方向追击廓尔喀，…}乾隆五十四年（1789年），清军越过喜马拉雅方向追击廓尔喀，边境协商与军事行动并行。核心取证：《清高宗实录》乾隆五十四年相关条；军机处档、战事方略及地方奏报。编年证词定位国家叙述中的时间，却需同地方或对方材料校验范围。来源保留：以乾隆五十四年（1789年）的同期文书为定位起点；官修记录、奏报或外交文本服务特定机构立场，具体日期、规模、地方执行与对手视角须以独立材料复核。
\subsection{1790—1799}
\eventanchor{EQ-1790-IMPERIAL-REGULATIONS-TIBET}{清廷制定整顿西藏事务的制度性安排…}乾隆五十五年（1790年），清廷制定整顿西藏事务的制度性安排，强化驻藏大臣和金瓶掣签等机制。核心取证：《清高宗实录》乾隆五十五年相关条；军机处档、理藩院档或相关方略。编年证词定位国家叙述中的时间，却需同地方或对方材料校验范围。来源保留：以乾隆五十五年（1790年）的同期文书为定位起点；官修记录、奏报或外交文本服务特定机构立场，具体日期、规模、地方执行与对手视角须以独立材料复核。
\eventanchor{EQ-1791-WHITE-LOTUS-PRECURSORS}{川楚陕基层社会的教门网络和地方压…}乾隆五十六年（1791年），川楚陕基层社会的教门网络和地方压力成为白莲教战争的背景。核心取证：《清高宗实录》乾隆五十六年相关条；地方志、督抚奏折及地方档案。编年证词定位国家叙述中的时间，却需同地方或对方材料校验范围。来源保留：以乾隆五十六年（1791年）的同期文书为定位起点；官修记录、奏报或外交文本服务特定机构立场，具体日期、规模、地方执行与对手视角须以独立材料复核。
\eventanchor{EQ-1792-GURKHA-SETTLEMENT}{清廷与廓尔喀战争结束后处理使节、…}乾隆五十七年（1792年），清廷与廓尔喀战争结束后处理使节、边防和西藏军政善后。核心取证：《清高宗实录》乾隆五十七年相关条；《筹办夷务始末》相关年卷与中外照会、条约文本。此类文件确认机构作出或收到的行政动作，不能跳过地方执行与反馈。来源保留：以乾隆五十七年（1792年）的同期文书为定位起点；官修记录、奏报或外交文本服务特定机构立场，具体日期、规模、地方执行与对手视角须以独立材料复核。
\eventanchor{EQ-1793-MACARTNEY-MISSION}{马戛尔尼使团抵达清廷，围绕贸易、…}乾隆五十八年（1793年），马戛尔尼使团抵达清廷，围绕贸易、礼仪和外交关系交涉。核心取证：《清高宗实录》乾隆五十八年相关条；《筹办夷务始末》相关年卷与中外照会、条约文本。此类文件确认机构作出或收到的行政动作，不能跳过地方执行与反馈。来源保留：以乾隆五十八年（1793年）的同期文书为定位起点；官修记录、奏报或外交文本服务特定机构立场，具体日期、规模、地方执行与对手视角须以独立材料复核。
\eventanchor{EQ-1794-WHITE-LOTUS-OUTBREAK}{白莲教相关起事在川楚陕地区扩展，…}乾隆五十九年（1794年），白莲教相关起事在川楚陕地区扩展，长期战争开始。核心取证：《清高宗实录》乾隆五十九年相关条；军机处档、战事方略及地方奏报。编年证词定位国家叙述中的时间，却需同地方或对方材料校验范围。来源保留：以乾隆五十九年（1794年）的同期文书为定位起点；官修记录、奏报或外交文本服务特定机构立场，具体日期、规模、地方执行与对手视角须以独立材料复核。
\eventanchor{EQ-1795-QIANLONG-ABDICATION}{乾隆帝禅位于嘉庆帝，自为太上皇，…}乾隆六十年（1795年），乾隆帝禅位于嘉庆帝，自为太上皇，仍保留重要政治影响。核心取证：《清高宗实录》乾隆六十年相关条；宫中朱批奏折、内阁大库题本与《清史稿》相关志传。此类文件确认机构作出或收到的行政动作，不能跳过地方执行与反馈。来源保留：以乾隆六十年（1795年）的同期文书为定位起点；官修记录、奏报或外交文本服务特定机构立场，具体日期、规模、地方执行与对手视角须以独立材料复核。
